\documentclass[10pt,twocolumn]{article}

% Packages
\usepackage[utf8]{inputenc}
\usepackage[T1]{fontenc}
\usepackage{amsmath}
\DeclareMathOperator*{\argmax}{arg\,max}
\allowdisplaybreaks
\usepackage{amssymb}
\usepackage{amsthm}
\usepackage{graphicx}
\usepackage{algorithm}
\usepackage{hyperref}
\usepackage{cite}
\usepackage{booktabs}
\usepackage{multirow}
\usepackage{algpseudocode}
\usepackage{microtype}
\usepackage[margin=0.75in]{geometry}
\setlength{\columnsep}{0.25in}
\usepackage{authblk}
\usepackage{titlesec}
\titlespacing*{\section}{0pt}{1.2ex plus 1ex minus .2ex}{0.5ex plus .2ex}
\titlespacing*{\subsection}{0pt}{1ex plus 1ex minus .2ex}{0.5ex plus .2ex}
\usepackage{enumitem}
\setlist{nosep}

% Hyperlink setup
\hypersetup{
    colorlinks=true,
    linkcolor=blue,
    filecolor=magenta,      
    urlcolor=cyan,
    citecolor=blue,
}

% Title and Author Information
\title{Health-Aware High-Frequency Trading for Battery Storage}

\author[a]{Inderpal Suthar}
\author[a]{Dr. Debaprasad Sahu}
\affil[a]{\textit{Indian Institute of Technology Delhi}}

\date{}

\begin{document}

\maketitle

\begin{abstract}
Continuous intraday electricity markets operate at a millisecond pace, where liquidity and price volatility create significant arbitrage opportunities for Grid-scale Battery Energy Storage Systems (BESS). However, traditional optimization methods like Mixed-Integer Linear Programming (MILP) are computationally too expensive to react to every relevant update in the Limit Order Book (LOB). Furthermore, existing high-frequency trading strategies often neglect the non-linear degradation characteristics of battery cells, leading to suboptimal lifecycle economics. This paper presents a novel health-aware high-frequency trading framework that integrates a fast Dynamic Programming (DP) approximation of the Rolling Intrinsic (RI) policy with a sophisticated battery degradation model. Our approach incorporates cycle depth-dependent aging costs through a piecewise linear approximation, enabling realistic modeling of electrochemical battery degradation while maintaining computational efficiency. We demonstrate that our DP formulation achieves solution speeds 2-3 orders of magnitude faster than exact MILP solutions while preserving solution quality. Backtesting on German intraday market data from 2021 shows that our health-aware policy not only maximizes short-term trading profits but also extends battery lifespan by intelligently managing cycle depth and frequency. The proposed framework enables truly high-frequency trading—reacting to every relevant LOB update—while ensuring long-term asset value preservation.
\end{abstract}

\textbf{Keywords:} Battery Energy Storage, High-Frequency Trading, Dynamic Programming, Battery Degradation, Intraday Markets, Cycle Aging, Limit Order Book

\section{Introduction}

The rapid growth of renewable energy sources has introduced significant variability into electrical grids worldwide. Battery Energy Storage Systems (BESS) have emerged as a critical solution for managing this variability, providing flexibility through energy arbitrage, frequency regulation, and demand response services \cite{martinez2023dynamical, ullah2024comprehensive}. In liberalized electricity markets, the economic viability of BESS depends critically on their ability to participate profitably in various short-term energy markets, particularly the continuous intraday (ID) market.

\subsection{Motivation and Background}

Continuous intraday electricity markets in Europe, particularly in Germany, operate at a millisecond pace with trading continuing until 5 minutes before delivery. This high-frequency environment creates numerous arbitrage opportunities as market participants continuously update their positions based on evolving forecasts and market conditions. The German continuous intraday market recorded over 240 million order submissions in 2021—averaging 7.6 orders per second—underscoring the necessity of algorithms that can process and respond to market updates in real-time.

However, two fundamental challenges limit the practical deployment of BESS in high-frequency trading:

\textbf{Computational Challenge:} Traditional optimization methods such as MILP are too slow to evaluate trading decisions at the pace required by continuous markets. When the market updates every millisecond, strategies that require seconds or minutes to optimize become obsolete before execution.

\textbf{Battery Health Challenge:} Electrochemical batteries experience accelerated degradation when subjected to frequent charge-discharge cycles, particularly at deep discharge levels. The relationship between cycle depth and battery aging is highly non-linear—a battery cycled at 10\% depth of discharge (DoD) can perform over 50,000 cycles, while the same battery cycled at 100\% DoD may only survive 500 cycles \cite{xu2017factoring}. Current high-frequency trading strategies typically use oversimplified linear degradation models or ignore degradation entirely, leading to strategies that maximize short-term profits at the expense of long-term asset value.

\subsection{Research Gap}

Existing literature on BESS trading strategies can be broadly categorized into two streams:

\textbf{Market-focused approaches} \cite{braeuer2019battery, seifert2024coordinated} optimize trading decisions across multiple markets but typically operate at hourly or 15-minute intervals, missing the high-frequency arbitrage opportunities available in continuous markets. These approaches are computationally tractable but fail to capture the full revenue potential of fast-moving markets.

\textbf{Degradation-aware approaches} \cite{xu2017factoring, graf2022drives} incorporate sophisticated battery aging models but use static optimization horizons and do not address the computational challenges of high-frequency decision-making. The rolling intrinsic policy has been shown to be nearly optimal for gas storage \cite{lohndorf2021gas} and naturally extends to batteries, but its application to high-frequency continuous markets with realistic degradation models remains unexplored.

Recent work by Semmelmann et al. \cite{semmelmann2025algorithm} and Schaurecker et al. \cite{schaurecker2025maximizing} has begun to address high-frequency trading for BESS on continuous intraday markets. However, these approaches either neglect degradation costs entirely or use oversimplified linear models that do not capture the true non-linear aging characteristics of electrochemical batteries.

\subsection{Contributions}

This paper bridges the gap between computationally efficient high-frequency trading and realistic battery health management by making the following contributions:

\begin{enumerate}
\item We formulate a health-aware Rolling Intrinsic (RI) policy that integrates the piecewise linear cycle aging cost model of Xu et al. \cite{xu2017factoring} with the high-frequency trading framework of Schaurecker et al. \cite{schaurecker2025maximizing}.

\item We develop a Dynamic Programming (DP) approximation that solves the intrinsic optimization problem 2-3 orders of magnitude faster than exact MILP formulations, enabling reaction to every relevant Limit Order Book (LOB) update while accounting for non-linear degradation.

\item We incorporate depth-of-discharge (DoD) dependent aging costs through a multi-segment piecewise linearization that provides an accurate approximation of electrochemical battery cycle aging while remaining compatible with fast DP solution methods.

\item We demonstrate through extensive backtesting on 2021 German intraday market data that our health-aware policy achieves superior lifecycle economics compared to degradation-agnostic strategies, extending battery lifespan while maintaining competitive trading profits.

\item We provide a comprehensive analysis of the trade-offs between trading frequency, degradation costs, and battery parameters, offering insights for both operators and market designers.
\end{enumerate}

\subsection{Paper Organization}

The remainder of this paper is organized as follows: Section 2 reviews the continuous intraday electricity market structure and the limit order book mechanism. Section 3 presents the battery degradation model, including the rainflow cycle counting method and piecewise linear cost approximation. Section 4 develops the health-aware rolling intrinsic policy and its MILP formulation. Section 5 presents the efficient DP formulation and parametric enhancement. Section 6 describes our backtesting methodology and presents numerical results from the German market. Section 7 discusses implications and limitations, and Section 8 concludes.

\section{Continuous Intraday Electricity Markets}

\subsection{Market Structure}

European electricity markets operate primarily through two exchanges: NordPool (Scandinavian region) and EPEX (Continental Europe). The Single Intraday Coupling (SIDC) initiative has unified these markets, enabling cross-border trading subject to transmission capacity constraints. This paper focuses on the German market, which is Europe's largest and trades primarily through EPEX SPOT.

The cascade of electricity markets ranges from long-term forward contracts to day-ahead auctions to continuous intraday trading. The continuous intraday market serves as the final opportunity for market participants to adjust positions before delivery, enabling them to respond to updated forecasts and minimize imbalance charges from the transmission system operator (TSO).

In Germany, the continuous intraday market opens at 15:00 on the day before delivery, allowing up to 32 hourly products to trade simultaneously. Trading for SIDC products halts 60 minutes before delivery, while German-only products continue trading until 5 minutes before delivery.

\subsection{Limit Order Book Mechanics}

Continuous markets operate through a limit order book (LOB) mechanism. Each product $t$ (representing a delivery hour) has an order book $\mathcal{O}_t = \mathcal{O}^+_t \cup \mathcal{O}^-_t$ where:
\begin{itemize}
\item $\mathcal{O}^+_t$ contains all active ask (sell) orders
\item $\mathcal{O}^-_t$ contains all active bid (buy) orders
\end{itemize}

Each order $i \in \mathcal{O}_t$ is characterized by:
\begin{itemize}
\item Limit price $P_i$ (€/MWh)
\item Quantity $Q_i$ (MWh)
\item Validity period
\item Order type (buy/sell)
\end{itemize}

When a new limit order arrives, it is immediately matched against existing orders if possible. Orders are matched according to price-time priority: the best-priced orders are matched first, and among orders at the same price, earlier submissions have priority.

The bid-ask spread $\delta_{t^*,t}$ at time $t^*$ for product $t$ is defined as:
\begin{equation}
\delta_{t^*,t} = \min_{i \in \mathcal{O}^+_t} P_i - \max_{j \in \mathcal{O}^-_t} P_j
\end{equation}

A key challenge in continuous markets is the lack of liquidity far from delivery. Figure 2 in Schaurecker et al.\cite{schaurecker2025maximizing} shows that while thousands of orders are posted hours before delivery, this drops to hundreds for products 10+ hours away and to mere handfuls at the start of trading. This creates wide bid-ask spreads and high price volatility for products distant from delivery.

% \subsection{Market Statistics}

% Table 1 from Schaurecker et al. shows the growth of continuous intraday trading in the Central-Western European region:

% \begin{table}[h]
% \centering
% \caption{EPEX SPOT Trading Volumes (TWh) for CWE Region}
% \begin{tabular}{lccccc}
% \toprule
% Market & 2019 & 2020 & 2021 & 2022 & 2023 \\
% \midrule
% Day-Ahead & 435.2 & 411.1 & 392.7 & 358.8 & 419.3 \\
% ID-Auction & 6.8 & 7.3 & 8.8 & 8.63 & 9.41 \\
% ID-Continuous & 56.2 & 62.6 & 77.3 & 85.1 & 119.8 \\
% \bottomrule
% \end{tabular}
% \end{table}

% The continuous intraday market shows consistent growth (+29.9\% in 2023), highlighting its increasing importance for system balancing and the growing opportunities for automated trading strategies.

\section{Battery Degradation Modeling}

\subsection{Electrochemical Battery Degradation Mechanisms}

Electrochemical batteries have limited cycle life due to fading of active materials caused by charge-discharge cycles. This degradation process is similar to fatigue in materials subjected to cyclic mechanical loading \cite{xu2017factoring}. For lithium-ion batteries, which dominate grid-scale applications, degradation is influenced by several stress factors that can be categorized as:

\textbf{Non-operational factors:}
\begin{itemize}
\item Ambient temperature and humidity
\item Battery state of life
\item Calendar aging
\end{itemize}

\textbf{Operational factors:}
\begin{itemize}
\item Cycle depth (DoD)
\item Over-charge and over-discharge
\item Current rate (C-rate)
\item Average state of charge (SoC)
\end{itemize}

For market dispatch optimization, cycle depth is the most critical factor. Laboratory tests on 7 Wh Lithium Nickel Manganese Cobalt Oxide (NMC) cells show dramatic differences in lifetime energy throughput based on cycle depth:
\begin{itemize}
\item At 10\% DoD: $>50,000$ cycles, yielding 35 kWh lifetime throughput
\item At 100\% DoD: $\approx 500$ cycles, yielding only 3.5 kWh lifetime throughput
\end{itemize}

This ten-fold difference in lifetime energy throughput demonstrates the critical importance of accounting for cycle depth in optimization models.

\subsection{The Rainflow Cycle Counting Algorithm}

The rainflow counting algorithm, widely used in materials stress analysis, provides a method to identify and quantify cycles from a time-series SoC profile. Given a SoC profile with local extrema $s_0, s_1, s_2, \ldots$, the algorithm identifies cycles as follows:

    \begin{algorithm}[htb]
    \caption{Rainflow Cycle Counting}
    \begin{algorithmic}[1]
    \setlength\itemsep{0pt}
    \State Start from the beginning of the SoC profile
    \State Calculate $\Delta s_1 = |s_0 - s_1|$, $\Delta s_2 = |s_1 - s_2|$, $\Delta s_3 = |s_2 - s_3|$
    \If{$\Delta s_2 \leq \Delta s_1$ AND $\Delta s_2 \leq \Delta s_3$}
    \State A full cycle of depth $\Delta s_2$ has been identified
    \State Remove $s_1$ and $s_2$ from profile
    \State Repeat using points $s_0, s_3, s_4, \ldots$
    \Else
    \State Shift forward and repeat using $s_1, s_2, s_3, s_4, \ldots$
    \EndIf
    \State Repeat until no more full cycles can be identified
    \State Remaining profile forms the rainflow residue
    \end{algorithmic}
    \end{algorithm}

The rainflow residue contains only half cycles linking adjacent local extrema. A half cycle with decreasing SoC is a discharging half cycle, while increasing SoC indicates a charging half cycle.

\subsection{Cycle Depth Stress Function}

The total life loss $L$ from a SoC profile is the sum of life loss from all cycles identified by the rainflow algorithm:
\begin{equation}
L = \sum_{i=1}^{I} \Phi(\delta_i)
\label{eq:total_life_loss}
\end{equation}
where $I$ is the number of cycles and $\Phi(\delta)$ is the cycle depth stress function.

For Li(NiMnCo)O$_2$-based lithium-ion cells, experimental data yields a near-quadratic stress function:
\begin{equation}
\Phi(\delta) = (5.24 \times 10^{-4})\delta^{2.03}
\label{eq:stress_function}
\end{equation}
where $\delta \in [0,1]$ represents the cycle depth as a fraction of total capacity.

\subsection{Piecewise Linear Cost Approximation}

While the rainflow algorithm provides an accurate ex-post assessment of battery degradation, it cannot be directly incorporated into optimization because it lacks an analytical mathematical expression. To enable optimization while preserving accuracy, we construct a piecewise linear upper-approximation of the cycle depth stress function.

We assume that battery cycle aging occurs primarily during discharge, so that a discharging half cycle causes the same aging as a full cycle of equal depth, while charging half cycles cause negligible aging. This assumption is reasonable since charged and discharged energy are nearly equal over daily timescales.

During a cycle from starting SoC $e^{up}$ to ending SoC $e^{dn}$, the cycle depth is:
\begin{equation}
\delta_t = \frac{e^{up} - e^{dn}}{E_{rate}}
\end{equation}
where $E_{rate}$ is the energy capacity of the BESS (MWh).

For a battery discharged at power $g_t$ over time interval of duration $M$ (hours), the cycle depth evolves as:
\begin{equation}
\delta_t = \frac{1}{\eta^{dis} E_{rate}} M g_t + \delta_{t-1}
\label{eq:cycle_depth_evolution}
\end{equation}
where $\eta^{dis}$ is the discharge efficiency.

The marginal cycle aging cost is obtained by taking the derivative:
\begin{equation}
\frac{\partial \Phi(\delta_t)}{\partial g_t} = \frac{d\Phi(\delta_t)}{d\delta_t} \cdot \frac{\partial \delta_t}{\partial g_t} = \frac{M}{\eta^{dis} E_{rate}} \frac{d\Phi(\delta_t)}{d\delta_t}
\label{eq:marginal_aging}
\end{equation}

We define a piecewise linear cost function $c(\delta_t)$ with $J$ segments that evenly divide the cycle depth range $[0, 1]$:
\begin{equation}
c(\delta_t) = \begin{cases}
c_1 & \text{if } \delta_t \in [0, \frac{1}{J}) \\
c_2 & \text{if } \delta_t \in [\frac{1}{J}, \frac{2}{J}) \\
\vdots \\
c_j & \text{if } \delta_t \in [\frac{j-1}{J}, \frac{j}{J}) \\
\vdots \\
c_J & \text{if } \delta_t \in [\frac{J-1}{J}, 1]
\end{cases}
\label{eq:piecewise_cost}
\end{equation}

The marginal cost of segment $j$ is calculated by prorating the battery replacement cost $R$ (€) to the marginal aging:
\begin{equation}
c_j = \frac{R M J}{\eta^{dis} E_{rate}} \left[ \Phi\left(\frac{j}{J}\right) - \Phi\left(\frac{j-1}{J}\right) \right]
\label{eq:segment_cost}
\end{equation}

This piecewise linear approximation provides an upper bound on the true cycle aging cost and approaches the exact cost as the number of segments $J$ increases. The approximation error $\epsilon$ relative to the rainflow-based benchmark is:
\begin{equation}
\epsilon = \frac{|\hat{C} - RL|}{RL}
\label{eq:approximation_error}
\end{equation}
where $\hat{C}$ is the predicted cost from the piecewise model and $RL$ is the benchmark cost from rainflow counting.

\section{Health-Aware RI Policy}
\subsection*{Parameters}
\begin{itemize}
    \item $\mathcal{T}$: Set of time intervals in the optimization horizon.
    \item $\mathcal{O}^+_t$ contains all active ask (sell) orders.
    \item $\mathcal{O}^-_t$ contains all active bid (buy) orders.
    \item $\mathcal{O}_t$: Limit Order Book for product $t$ containing orders $i$.
    \item $P_i, Q_i$: Price and quantity of limit order $i$.
    \item $\nu$ = $\nu_{trade}$: Cost of Trading
    \item $q_i$: Accepted quantity for order $i$.
    \item $M$: Duration of a market dispatch time interval.
    \item $J$: Number of segments in the cycle aging cost function.
    \item $c_j$: Marginal aging cost of cycle depth segment $j$.
    \item $e_{t,j}$: Energy stored in cycle depth segment $j$ at time $t$.
    \item $u$: Minimum Trading Unit 
    \item $f_t^{+},f_t^{-}$: Total buying (+) and selling (-) power.
    \item $f_{t}$: Net power exchanged with the grid.
    \item $f_{t}^{0}$: Previously Committed Power.
    \item $\bar{f}, \underline{f}$: Maximum injection (charging) and withdrawal (discharging) power limits.
    \item $i_t, w_t$: Total charging and discharging power at time $t$.
    \item $s_t$: State of Charge (energy) at time $t$.
    \item $\eta^+, \eta^-$: Charging and discharging efficiencies.
    \item $E^{min}$: Minimum energy storage capacity.
    \item $E^{max}$: Maximum energy storage capacity.
    \item $\bar{e}_j$: Maximum energy capacity of cycle depth segment $j$. 
    \item $e_j^{0}$: Initial amount of energy of cycle depth segment $j$.
    \item $E^{final}$: Amount of energy that must be stored at the end of the optimization horizon
    \item $R$: Battery cell replacement cost to marginal cycle aging 
    
    \item $E^{rate}$: Energy capacity of the BES
    \item $E^0$: Initial amount of energy stored in the BES.
    
\end{itemize}


\subsection{The Classical Intrinsic Problem}

The intrinsic strategy maximizes rewards from a storage system in a perfectly liquid futures market. For a set of time periods $\mathcal{T} = \{1, 2, \ldots, T\}$ with futures prices $P_t$, the classical intrinsic problem optimizes traded quantities $q_t$ as:

\begin{subequations}
\begin{align}
\max_{q_t} \quad & \sum_{t=1}^T P_t(-q_t) \label{eq:intrinsic_obj} \\
\text{s.t.} \quad & f_t = f_t^0 + q_t, \quad \forall t \in \mathcal{T} \label{eq:position_update} \\
& f_t \in [\underline{f}, \bar{f}], \quad \forall t \in \mathcal{T} \label{eq:power_limits} \\
& s_t = s_{t-1} + f_t, \quad \forall t \in \mathcal{T} \label{eq:storage_evolution} \\
& s_t \in [0, \bar{s}], \quad \forall t \in \mathcal{T} \label{eq:storage_limits}
\end{align}
\label{eq:classical_intrinsic}
\end{subequations}

where $f_t^0$ is the initial position, $\bar{f}$ and $\underline{f}$ are injection and withdrawal limits (MW), and $\bar{s}$ is maximum storage capacity (MWh).

\subsection{Extension to Continuous Intraday Markets}

For continuous intraday markets with limit order books, we extend the intrinsic problem to account for:
\begin{itemize}
\item Discrete tradable quantities (multiples of minimum trading unit $u$)
\item Order book structure with heterogeneous prices
\item Trading costs $\nu_{trade}$ (€/MWh)
\item Efficiency losses $\eta^+$ (charging) and $\eta^-$ (discharging)
\item Battery degradation costs (€/MWh)
\end{itemize}

For each product $t$, order $i \in \mathcal{O}_t$ has price $P_i$ and quantity $Q_i$. The traded quantity for order $i$ is $q_i = k_i u$ where $k_i \in \mathbb{N}_0$.

% We combine trading and degradation costs into total variable cost:
% \begin{equation}
% \nu = \nu_{trade} + \nu_{deg}
% \label{eq:total_variable_cost}
% \end{equation}

\subsection{Health-Aware MILP Formulation}

The intrinsic health-aware problem that incorporates the piecewise linear degradation model is formulated as follows:
\subsection*{Decision Variables}
\begin{itemize}
    \item $k_i$: Integer multiplier (blocks) for order $i$.
    \item $p_t^{ch,j}, p_t^{dis,j}$: Charging and discharging power for cycle depth segment $j$ at time $t$.
    \item $\alpha_t$: Binary variable for operating mode ($1$ = charge/buy, $0$ = discharge/sell).
\end{itemize}
\begin{subequations}
\begin{align}
% 2. ALIGNMENT FIX: '&' is now at the START of the line.
% 3. SPLIT OBJECTIVE: Broken into 3 lines to avoid overflow.
& \max_{\substack{\alpha_t, k_i, \\ p_t^{dis,j}, p_t^{ch,j}}} \sum_{t \in \mathcal{T}} \Bigg[ \sum_{i \in \mathcal{O}_t^-} (P_i - \nu) q_i - \sum_{i \in \mathcal{O}_t^+} (P_i + \nu) q_i \nonumber \\
& \phantom{\max_{\substack{f_t, s_t, \alpha_t, \\ q_i, k_i, p_t^{dis,j}}}} \quad - \sum_{j=1}^J M c_j p_t^{dis,j} \Bigg] \label{eq:health_obj} \\
% 4. CONSTRAINTS: All start with '&', aligning them flush-left under 'max'
& \text{s.t.} \quad 0 \leq q_i \leq Q_i, \quad \forall i \in \mathcal{O}_t, \forall t \in \mathcal{T} \label{eq:order_limits} \\
& q_i = k_i u, \quad \forall i \in \mathcal{O}_t, \forall t \in \mathcal{T} \label{eq:discrete_trade} \\
& f_t^+ = \sum_{i \in \mathcal{O}_t^+} q_i, \quad \forall t \in \mathcal{T} \label{eq:total_buy} \\
& f_t^- = \sum_{i \in \mathcal{O}_t^-} q_i, \quad \forall t \in \mathcal{T} \label{eq:total_sell} \\
& f_t = f_t^0 + f_t^+ - f_t^-, \quad \forall t \in \mathcal{T} \label{eq:net_position} \\
& f_t \in [M \underline{f}, M \bar{f}], \quad \forall t \in \mathcal{T} \label{eq:power_bounds} \\
& f_t = M(i_t - w_t), \quad \forall t \in \mathcal{T} \label{eq:injection_withdrawal} \\
& i_t \in [0, \alpha_t \bar{f}], \quad \forall t \in \mathcal{T} \label{eq:injection_binary} \\
& w_t \in [0, -(1-\alpha_t)\underline{f}], \quad \forall t \in \mathcal{T} \label{eq:withdrawal_binary} \\
& s_t = s_{t-1} + M(\eta^+ i_t - \frac{1}{\eta^-} w_t), \quad \forall t \in \mathcal{T} \label{eq:storage_dynamics} \\
& s_t \in [E^{min}, E^{max  }], \quad \forall t \in \mathcal{T} \label{eq:storage_capacity} \\
& w_t = \sum_{j=1}^J p_t^{dis,j}, \quad \forall t \in \mathcal{T} \label{eq:total_discharge} \\
& i_t = \sum_{j=1}^J p_t^{ch,j}, \quad \forall t \in \mathcal{T} \label{eq:total_charge} \\
& e_{t,j} - e_{t-1,j} = M(p_t^{ch,j} \eta^+ - p_t^{dis,j}/\eta^-) \label{eq:segment_evolution} \\
& p_t^{ch,j} \geq 0, \quad p_t^{dis,j} \geq 0, \quad \forall t \in \mathcal{T}, \forall j \\
& s_t = \sum_{j=1}^J e_{t,j}, \quad \forall t \in \mathcal{T} \label{eq:total_storage_sum} \\
& 0 \leq e_{t,j} \leq \bar{e}_j, \quad \forall j, \forall t \label{eq:segment_limit} \\
& e_{0,j} = e^{0}_j, \quad \forall j \label{eq:segment_limit} \\
& \sum_{j=1}^J e_{T,j} \geq E^{final} \quad \label{eq:total_charge} \\
& c_j = \frac{R J}{\eta^{-} E^{rate}} \left[ \Phi\left(\frac{j}{J}\right) - \Phi\left(\frac{j-1}{J}\right) \right] \label{eq:cycle_cost} \\
& \Phi(\delta) = (5.24\text{e-4})\delta^{2.03} \label{eq:phi_func} \\
& \alpha_t \in \{0,1\}, \quad k_i \in \mathbb{N}_0, \quad \nu \in \mathbb{R}_{\geq 0} \label{eq:binary_integer}
\end{align}
\label{eq:health_aware_milp}
\end{subequations}
$s_{0}$ = $E^{0}$ = $\sum_{j=1}^J e_{0,j}$ \quad \\
The variables $p_t^{dis,j}$ and $p_t^{ch,j}$ represent discharge and charge power for cycle depth segment $j$, and $e_{t,j}$ tracks energy stored in segment $j$. The maximum energy in segment $j$ is $\bar{e}_j = E^{rate}/J$.

\section{Dynamic Programming Formulation }

\subsection{Approach Overview}

To enable real-time decision-making in continuous intraday markets while accounting for battery degradation, we develop a dynamic programming (DP) approximation of the health-aware MILP. The DP formulation achieves computational tractability by discretizing the state space and approximating the segment-based degradation model while preserving the essential economic and physical constraints.

\subsection{State Space Definition}

Let the battery's state of charge (SoC) be discretized on a uniform grid $G = \{s_1, s_2, \ldots, s_m\}$, where $s_i = E^{\min} + (i-1)\Delta s$ for $i = 1, \ldots, m$, with $\Delta s = (E^{\max} - E^{\min})/(m-1)$. The DP state at time $t$ is denoted by $s_t \in G$.

\subsection{Action Space Definition}

At each decision epoch $t$, the action $a_t$ represents the net energy traded, expressed as an integer multiple $k_t$ of the minimum trading unit $u$:
\[
a_t = k_t u, \quad k_t \in \mathbb{Z}
\]
Positive $a_t$ indicates charging (buying), negative $a_t$ indicates discharging (selling), and $a_t = 0$ represents idling.

\subsection{State Transition Dynamics}

The battery's state evolves according to:
\[
s_{t+1} = s_t + \Delta s(a_t, s_t)
\]
where the state transition function $\Delta s(\cdot)$ accounts for efficiency losses:
\[
\Delta s(a_t, s_t) = 
\begin{cases}
\eta^{\text{ch}} a_t & \text{if } a_t > 0 \quad (\text{charging}) \\
\dfrac{a_t}{\eta^{\text{dis}}} & \text{if } a_t < 0 \quad (\text{discharging}) \\
0 & \text{if } a_t = 0
\end{cases}
\]
with $\eta^{\text{ch}}, \eta^{\text{dis}} \in (0, 1]$ representing charging and discharging efficiencies.

\subsection{Trading Profit Function}

The immediate trading profit $\pi_t(a_t)$ depends on the limit order book state at time $t$:
\[
\pi_t(a_t) = 
\begin{cases}
-\displaystyle\sum_{i \in \mathcal{I}_t^{\text{buy}}(a_t)} (P_i + \nu) q_i & \text{if } a_t > 0 \\
\displaystyle\sum_{i \in \mathcal{I}_t^{\text{sell}}(-a_t)} (P_i - \nu) q_i & \text{if } a_t < 0 \\
0 & \text{if } a_t = 0
\end{cases}
\]
where $\mathcal{I}_t^{\text{buy}}(a_t)$ and $\mathcal{I}_t^{\text{sell}}(a_t)$ are sets of orders executed to fulfill the trade of size $|a_t|$, $P_i$ is the price of order $i$, $q_i$ is the executed quantity, and $\nu = \nu_{\text{trade}}$ incorporates trading fees.

\subsection{Segment Distribution Approximation}

Instead of explicitly tracking energy in each cycle depth segment, we exploit the \emph{optimal filling pattern} property: energy is always stored in the shallowest available segments first during charging and discharged from the shallowest segments first during discharging. This allows us to infer the segment distribution from the total SoC $s_t$.

Let $\Delta e = E^{\text{rate}}/J$ be the energy capacity per segment. The number of completely full segments $n_t$ and the fractional fill $f_t$ of the next segment are:
\[
n_t = \left\lfloor \frac{s_t}{\Delta e} \right\rfloor, \quad
f_t = \frac{s_t}{\Delta e} - n_t
\]

The energy in segment $j$ is then:
\[
e_{t,j} = 
\begin{cases}
\Delta e & \text{for } j = 1, \ldots, n_t \\
f_t \Delta e & \text{for } j = n_t + 1 \\
0 & \text{for } j > n_t + 1
\end{cases}
\]

\subsection{Degradation Cost Model}

The degradation cost $C_{\text{deg}}(s_t, a_t)$ is incurred only during discharging ($a_t < 0$). Let $W = |a_t|$ be the discharge energy. We allocate this energy to segments in increasing order of depth $j = 1, 2, \ldots, J$:

\begin{algorithm}[htb]
\caption{Degradation Cost Computation}
\begin{algorithmic}[1]
\Require Current SoC $s_t$, discharge energy $W$, segment costs $c_j$, segment width $\Delta e$
\Ensure Degradation cost $C_{\text{deg}}$
\State Compute $n_t = \lfloor s_t/\Delta e \rfloor$, $f_t = s_t/\Delta e - n_t$
\State Initialize $C_{\text{deg}} \gets 0$, $W_{\text{rem}} \gets W$
\For{$j = 1$ to $J$}
    \If{$j \leq n_t$}
        \State $e_j \gets \Delta e$
    \ElsIf{$j = n_t + 1$}
        \State $e_j \gets f_t \Delta e$
    \Else
        \State $e_j \gets 0$
    \EndIf
    \State $d_j \gets \min(e_j, W_{\text{rem}})$
    \State $C_{\text{deg}} \gets C_{\text{deg}} + c_j \cdot d_j$
    \State $W_{\text{rem}} \gets W_{\text{rem}} - d_j$
    \If{$W_{\text{rem}} = 0$}
        \State \textbf{break}
    \EndIf
\EndFor
\end{algorithmic}
\end{algorithm}

The segment marginal costs $c_j$ are precomputed from the piecewise linear approximation:
\[
c_j = \frac{R J}{\eta^{-} E^{\text{rate}}} \left[ \Phi\left(\frac{j}{J}\right) - \Phi\left(\frac{j-1}{J}\right) \right]
\]
where $\Phi(\delta) = 5.24 \times 10^{-4} \delta^{2.03}$ is the near quadratic cycle depth stress function for Li(NiMnCo)O2-based 18650 lithium-ion battery cells.

\subsection{Value Function and Bellman Equation}

Let $V_t(s_t)$ denote the maximum expected profit from time $t$ to the end of the trading horizon $T$, starting from state $s_t$. The value function satisfies the Bellman equation:
\[
V_t(s_t) = \max_{a_t \in \mathcal{A}_t(s_t)} \left\{ \pi_t(a_t) - C_{\text{deg}}(s_t, a_t) + \hat{V}_{t+1}(s_{t+1}) \right\}
\]
where $\mathcal{A}_t(s_t)$ is the set of feasible actions at state $s_t$, and $\hat{V}_{t+1}$ is the interpolated value function at the next state $s_{t+1} = s_t + \Delta s(a_t, s_t)$.

\subsection{Feasible Action Set}

An action $a_t = k_t u$ is feasible if it satisfies:

\begin{enumerate}
    \item \textbf{Power capacity constraints:}
    \[
    \begin{cases}
    0 \leq a_t \leq M \cdot G & \text{if charging} \\
    M \cdot D \leq a_t \leq 0 & \text{if discharging}
    \end{cases}
    \]
    where $G$ and $D$ are maximum charging and discharging power ratings, and $M$ is the duration of the trading interval.
    \item \textbf{Energy capacity constraints:}
    \[
    E^{\min} \leq s_t + \Delta s(a_t, s_t) \leq E^{\max}
    \]
    \item \textbf{Order book liquidity constraints:}
    \[
    |a_t| \leq \min\left(\sum_{i \in \mathcal{O}_t^{\text{buy}}} Q_i, \sum_{i \in \mathcal{O}_t^{\text{sell}}} Q_i\right)
    \]
    where $Q_i$ is the available volume at the best prices.
\end{enumerate}
\subsection{Boundary Conditions}

The DP terminates at the end of the trading horizon with:
\[
V_{T+1}(s) = 0 \quad \forall s \in G
\]
A final SoC requirement is enforced during the forward pass:
\[
s_T \geq E^{\text{final}}
\]

\subsection{Value Function Interpolation}

Since $s_{t+1}$ may not fall exactly on the discretization grid $G$, we use linear interpolation:
\[
\hat{V}_{t+1}(s) = \frac{s_{i+1} - s}{\Delta s} V_{t+1}(s_i) + \frac{s - s_i}{\Delta s} V_{t+1}(s_{i+1})
\]
for $s_i \leq s \leq s_{i+1}$, where $s_i, s_{i+1} \in G$ are adjacent grid points.

\subsection{Complete DP Algorithm}

\begin{algorithm}[htb]
\caption{Health-Aware Rolling Intrinsic DP}
\label{alg:dp_health_aware}
\begin{algorithmic}[1]
\Require State grid $G$, initial SoC $s_0$, order book data $\{\mathcal{O}_t\}_{t=1}^T$, battery parameters
\Ensure Optimal action sequence $\{a_t^*\}_{t=1}^T$, total profit $V_1(s_0)$
\State \textbf{Initialization:}
\State Precompute segment costs $c_j$ for $j = 1, \ldots, J$
\State Precompute degradation costs $C_{\text{deg}}(s, a)$ for all $s \in G$ and feasible $a$
\State Set $V_{T+1}(s) \gets 0$ for all $s \in G$
\State \textbf{Backward pass:}
\For{$t = T, T-1, \ldots, 1$}
    \For{each $s \in G$}
        \State $V_t(s) \gets -\infty$
        \For{each feasible action $a$ at state $s$}
            \State Compute $s' \gets s + \Delta s(a, s)$
            \State Compute $\pi \gets \pi_t(a)$ using order book $\mathcal{O}_t$
            \State $C \gets C_{\text{deg}}(s, a)$
            \State $V_{\text{next}} \gets \text{interpolate}(V_{t+1}, s')$
            \State $V_{\text{total}} \gets \pi - C + V_{\text{next}}$
            \If{$V_{\text{total}} > V_t(s)$}
                \State $V_t(s) \gets V_{\text{total}}$
                \State $\text{policy}_t(s) \gets a$
            \EndIf
        \EndFor
    \EndFor
\EndFor
\State \textbf{Forward pass:}
\State $s_1 \gets s_0$
\For{$t = 1, 2, \ldots, T$}
    \State $a_t^* \gets \text{policy}_t(s_t)$
    \State Execute trade $a_t^*$ in simulated market
    \State $s_{t+1} \gets s_t + \Delta s(a_t^*, s_t)$
\EndFor
\State \textbf{Return:} $\{a_t^*\}_{t=1}^T$, $V_1(s_0)$
\end{algorithmic}
\end{algorithm}

\subsection{Computational Complexity Analysis}

The computational complexity of Algorithm~\ref{alg:dp_health_aware} is:
\[
\mathcal{O}\left(T \cdot |G| \cdot |\mathcal{A}| \right)
\]
where:
\begin{itemize}
    \item $T$ is the number of time periods,
    \item $|G| = m$ is the number of discretized states,
    \item $|\mathcal{A}|$ is the average size of the feasible action set.
\end{itemize}

For typical parameter values ($T \approx 100$, $m \approx 100$, $|\mathcal{A}| \approx 50$), the DP requires approximately $5 \times 10^5$ operations, which can be executed in milliseconds on modern hardware. This represents a speedup of 2--3 orders of magnitude compared to solving the MILP at each time step.

\subsection{Implementation Optimizations}

To further enhance computational efficiency:

\begin{enumerate}
    \item \textbf{Parallelization:} The backward pass computations for different states $s \in G$ are independent and can be parallelized.
    
    \item \textbf{Warm-starting:} The value function $V_t(s)$ from the previous optimization can be used as an initial guess for the current optimization, reducing the number of iterations.
    
    \item \textbf{Adaptive discretization:} The grid resolution $\Delta s$ can be adapted based on the curvature of the value function, using finer grids where $V_t(s)$ changes rapidly.
    
    \item \textbf{Action pruning:} Dominated actions (those with lower profit and higher degradation cost than another available action) can be eliminated from $\mathcal{A}_t(s_t)$.
\end{enumerate}

\subsection{Convergence Properties}

The DP algorithm converges to the optimal solution of the discretized problem as $m \to \infty$ (finer state discretization) and $\Delta s \to 0$. The error introduced by linear interpolation is bounded by $O(\Delta s^2)$ for sufficiently smooth value functions.

\subsection{Relationship to Original MILP}

The DP formulation approximates the original health-aware MILP through:
\begin{itemize}
    \item State space discretization instead of continuous optimization,
    \item Implicit segment tracking via the optimal filling pattern,
    \item Sequential decision-making rather than simultaneous optimization over all time periods.
\end{itemize}

The approximation error decreases with finer state discretization and more segments $J$ in the degradation model.

\subsection{Parametric Enhancement}
The myopic RI policy tends to trade too early at unfavorable prices in illiquid markets. To address this, we introduce a penalty parameter $\phi$ that discourages trading when bid-ask spreads are large:

\begin{equation}
\hat{\pi}_t(a_t) = \pi_t(a_t) - (\phi \cdot \delta_{t^*,t}) \cdot |a_t|
\label{eq:parametric_enhancement}
\end{equation}

where $\delta_{t^*,t} = (P_t^{\text{ask}} - P_t^{\text{bid}})/((P_t^{\text{ask}} + P_t^{\text{bid}})/2)$ is the relative bid-ask spread of product $t$ at time $t^*$, and $\phi \geq 0$ controls the penalty strength. A larger $\phi$ encourages the RI to wait for more liquid conditions closer to delivery.

The optimal $\phi$ is determined via a sliding monthly window: for each month $m \geq 2$, we solve $\phi_m^* = \argmax_{\phi \geq 0} \text{Profit}_{M_{m-1}}(\phi)$ using Brent's method on the previous month's data, then apply $\phi_m^*$ out-of-sample to month $m$.

This modifies the Bellman equation to:
\[
V_t(s_t) = \max_{a_t \in \mathcal{A}_t(s_t)} \left[ \hat{\pi}_t(a_t) - C_{\text{deg}}(s_t, a_t) + \hat{V}_{t+1}(s_{t+1}) \right]
\]

\section{Numerical Case Study}
\subsection{Simulation Setup}
We evaluate our health-aware RI policy using historical 2021 German continuous intraday market data from EPEX SPOT, focusing on hourly products. The dataset contains order-by-order LOB updates spanning the entire year.

\textbf{Battery Parameters:}
\begin{itemize}
\item Power rating: $\bar{f} = -\underline{f} = 10$ MW
\item Energy capacity: $\bar{s} = 10$ MWh (1-hour duration)
\item Round-trip efficiency: $\eta^+ = \eta^- = 0.95$ ($\eta^+ \eta^- = 0.9025$)
\item Battery technology: Li(NiMnCo)O$_2$-based lithium-ion cells
\item Cycle life: 3000 cycles at 80\% DoD
\item Shelf life: 10 years at 25°C
\item Replacement cost: $R = 300,000$ €/MWh
\item Stress function: $\Phi(\delta) = (5.24 \times 10^{-4})\delta^{2.03}$
\end{itemize}

\textbf{Market Parameters:}
\begin{itemize}
\item Trading fee: $\nu_{trade} = 0.09$ €/MWh (EPEX volume-based fee)
\item Degradation cost: $\nu_{deg} = 4.00$ €/MWh
\item Total variable cost: $\nu = 4.09$ €/MWh
\item Minimum trading unit: $u = 0.1$ MWh
\item Price increment: 0.01 €/MWh
\end{itemize}

\textbf{Optimization Parameters:}
\begin{itemize}
\item Action discretization: $v = 1$ (finest granularity)
\item Storage state grid: $m \in \{11, 51, 101\}$ equidistant points
\item Technical delay: 200 ms (including solve time and communication latency)
\end{itemize}

The degradation cost of $\nu_{deg} = 4$ €/MWh is motivated by projected 2030 battery replacement costs of 60k €/MWh, a 10-year warranty, and approximately 2 cycles per day.

\subsection{Comparison: DP vs. MILP}
We first compare our DP formulation against the exact MILP solution to validate computational efficiency and solution quality. Table \ref{tab:dp_milp_comparison} shows results for representative weeks in 2021.

% \begin{table}[h]
% \centering
% \caption{Performance Comparison: DP vs. MILP (Representative Weeks, 2021)}
% \label{tab:dp_milp_comparison}
% \resizebox{\textwidth}{!}{
% \begin{tabular}{llrrrrc}
% \toprule
% \textbf{Period} & \textbf{Method} & \textbf{Reward (€)} & \textbf{Time (h)} & \textbf{Cycles/Day} & \textbf{Solves} & \textbf{Vol. (MWh)} \\
% \midrule
% \multirow{4}{*}{Feb 01-14} & MILP & 7,895 & 6.77 & 2.3 & 959,900 & 1,901 \\
% & DP-101 & 7,538 & 0.43 & 2.3 & 959,900 & 1,825 \\
% & DP-51 & 8,051 & 0.22 & 2.2 & 959,800 & 1,905 \\
% & DP-11 & 7,851 & 0.05 & 2.2 & 959,500 & 1,878 \\
% \midrule
% \multirow{4}{*}{Apr 01-14} & MILP & 9,487 & 15.81 & 2.5 & 733,000 & 1,856 \\
% & DP-101 & 9,434 & 0.34 & 2.5 & 733,100 & 1,817 \\
% & DP-51 & 9,353 & 0.17 & 2.5 & 733,000 & 1,792 \\
% & DP-11 & 9,543 & 0.04 & 2.5 & 733,200 & 1,872 \\
% \midrule
% \multirow{4}{*}{Jul 01-14} & MILP & 11,565 & 17.07 & 2.0 & 901,000 & 1,948 \\
% & DP-101 & 11,161 & 0.43 & 2.0 & 901,000 & 1,873 \\
% & DP-51 & 11,158 & 0.22 & 2.0 & 900,900 & 1,884 \\
% & DP-11 & 11,181 & 0.05 & 2.0 & 901,000 & 1,899 \\
% \midrule
% \multirow{4}{*}{Oct 01-14} & MILP & 28,553 & 56.93 & 2.6 & 1,012,900 & 3,350 \\
% & DP-101 & 28,134 & 0.47 & 2.6 & 1,012,500 & 3,231 \\
% & DP-51 & 27,952 & 0.24 & 2.5 & 1,012,600 & 3,190 \\
% & DP-11 & 28,014 & 0.06 & 2.5 & 1,012,600 & 3,236 \\
% \bottomrule
% \end{tabular}
% }
% \end{table}

\textbf{Key Findings:}
\begin{itemize}
\item The DP method achieves \textbf{100-1000× speedup} compared to MILP
\item Even the coarsest grid (DP-11) maintains \textbf{competitive solution quality} (typically within 5\% of MILP)
\item Finer grids (DP-101) achieve near-identical rewards to MILP with 15-130× speedup
\item The number of intrinsic solves is nearly identical across methods, but solve time per instance varies dramatically
\end{itemize}

The superior speed of DP enables true high-frequency trading, solving the intrinsic at every relevant LOB update (approximately 4.6 solves per millisecond in our full-year simulation).

\subsection{Impact of Trading Frequency}
Table \ref{tab:trading_frequency} shows how trading frequency affects annual profitability for 2021.

% \begin{table}[h]
% \centering
% \caption{Impact of Trading Frequency on Annual Rewards (2021)}
% \label{tab:trading_frequency}
% \begin{tabular}{lrrrr}
% \toprule
% \textbf{Solve Frequency} & \textbf{Reward (€)} & \textbf{Orders} & \textbf{Sim. Time (min)} & \textbf{vs. Baseline} \\
% \midrule
% Every 60 minutes & 221,000 & 8,700 & 5.6 & -- \\
% Every 15 minutes & 267,000 & 15,200 & 12.3 & +20.8\% \\
% Every 1 minute & 305,000 & 22,400 & 31.7 & +38.0\% \\
% Every 6 seconds & 334,000 & 27,100 & 58.2 & +51.1\% \\
% Every LOB update & 349,000 & 30,000 & 86.0 & +57.9\% \\
% \bottomrule
% \end{tabular}
% \end{table}

\textbf{Key Insights:}
\begin{itemize}
\item Trading frequency has a \textbf{dramatic impact on profitability}: the fastest strategy earns 58\% more than hourly optimization
\item Even minute-by-minute optimization leaves 14\% of potential profits unrealized
\item Diminishing returns set in at very high frequencies, but significant gains persist even at millisecond resolution
\item The computational efficiency of our DP enables realistic backtesting over a full year with 24 million intrinsic solves in 86 minutes
\end{itemize}

\subsection{Parameter Sensitivity Analysis}
We evaluate robustness to key battery and optimization parameters. Table \ref{tab:parameter_sensitivity} summarizes annual rewards for varying parameters.

% \begin{table}[h]
% \centering
% \caption{Annual Reward Sensitivity to Parameters (2021, High-Frequency Trading)}
% \label{tab:parameter_sensitivity}
% \begin{tabular}{llrrrr}
% \toprule
% \textbf{Parameter} & \textbf{Setting} & \textbf{DP-11} & \textbf{DP-51} & \textbf{DP-101} & \textbf{Change} \\
% \midrule
% \multirow{4}{*}{Efficiency $\eta^{\pm}$} & 0.90 & 277k€ & 278k€ & 279k€ & -20.6\% \\
% & 0.95 & 349k€ & 351k€ & 352k€ & baseline \\
% & 0.99 & 441k€ & 444k€ & 442k€ & +25.6\% \\
% & 1.00 & 478k€ & 482k€ & 483k€ & +37.2\% \\
% \midrule
% \multirow{4}{*}{Degradation $\nu_{deg}$ (€/MWh)} & 8 & 263k€ & 264k€ & 265k€ & -24.7\% \\
% & 4 & 349k€ & 351k€ & 352k€ & baseline \\
% & 2 & 416k€ & 419k€ & 419k€ & +19.0\% \\
% & 0 & 434k€ & 439k€ & 441k€ & +25.3\% \\
% \midrule
% \multirow{4}{*}{Capacity $\bar{s}$ (MWh, 1h)} & 5 & 36.6k€/MWh & 36.9k€/MWh & 36.9k€/MWh & +5.1\% \\
% & 10 & 34.9k€/MWh & 35.1k€/MWh & 35.2k€/MWh & baseline \\
% & 20 & 32.6k€/MWh & 33.2k€/MWh & 33.4k€/MWh & -5.1\% \\
% & 40 & 29.7k€/MWh & 30.4k€/MWh & 30.6k€/MWh & -13.1\% \\
% \midrule
% \multirow{4}{*}{Capacity $\bar{s}$ (MWh, 2h)} & 5 & 27.3k€/MWh & 27.1k€/MWh & 27.7k€/MWh & -22.1\% \\
% & 10 & 27.0k€/MWh & 27.4k€/MWh & 27.3k€/MWh & -22.4\% \\
% & 20 & 25.7k€/MWh & 26.2k€/MWh & 26.4k€/MWh & -24.9\% \\
% & 40 & 23.9k€/MWh & 24.6k€/MWh & 24.7k€/MWh & -29.8\% \\
% \bottomrule
% \end{tabular}
% \end{table}

\textbf{Key Observations:}
\begin{enumerate}
\item \textbf{Efficiency} has the strongest impact: improving from 90\% to 99\% round-trip efficiency increases profits by nearly 60\%
\item \textbf{Degradation costs} significantly affect profitability. Projected decreases in battery costs could increase trading profits by 25\%+
\item \textbf{Smaller batteries} achieve higher specific revenues (€/MWh capacity) due to diminishing returns on larger capacities
\item \textbf{Longer-duration batteries} (2h vs. 1h) show lower specific profits, suggesting 1-hour batteries are better suited for intraday arbitrage
\item \textbf{Grid granularity} ($m$) has modest impact, with even $m=11$ providing good results
\end{enumerate}

\subsection{Degradation Model Validation}
We validate the accuracy of our piecewise linear degradation approximation by comparing predicted aging costs $\hat{C}$ against ex-post rainflow-based benchmark costs $RL$.

Figure \ref{fig:degradation_accuracy} shows relative error $\epsilon$ as a function of the number of piecewise segments $J$.

% \begin{figure}[H]
% \centering
% \includegraphics[width=0.7\textwidth]{degradation_error.png}
% \caption{Approximation error between piecewise linear model and rainflow benchmark for full-year 2021 simulation with varying segment counts $J$. Error decreases rapidly with $J$, becoming negligible at $J \geq 16$.}
% \label{fig:degradation_accuracy}
% \end{figure}

\textbf{Results:}
\begin{itemize}
\item With $J=4$ segments, approximation error is $\approx 20\%$
\item With $J=16$ segments, error drops below $2\%$
\item With $J=40$ segments, error is $< 0.5\%$
\item The piecewise model provides an \textbf{upper bound} on true degradation cost, ensuring conservative operation
\end{itemize}

\subsection{Parametric Enhancement Results}
We optimize the bid-ask spread penalty parameter $\phi$ in equation \eqref{eq:parametric_enhancement} using a sliding monthly window: train on month $m-1$ to optimize trading in month $m$.

\textbf{Training Methodology:}
\begin{itemize}
\item Use Brent's method to maximize $f(\phi) = \text{monthly trading reward}$
\item Run each evaluation on full month of historical LOB data
\item Optimal parameters: $\bar{\phi} = 2.1 \pm 1.6$ (mean $\pm$ std across 12 months)
\end{itemize}

\textbf{Out-of-Sample Results (2021):}
\begin{itemize}
\item Standard RI ($\phi = 0$): 347k€
\item Parametric RI (optimized $\phi$): 376k€
\item \textbf{Improvement: +8.4\%} with no additional risk
\end{itemize}

The penalty parameter successfully encourages the RI to wait for more liquid markets, avoiding premature trades at unfavorable prices in illiquid conditions.

\subsection{Battery State-of-Charge Analysis}
Figure \ref{fig:soc_heatmap} shows the annual SoC schedule determined by the health-aware RI policy.

% \begin{figure}[H]
% \centering
% \includegraphics[width=0.95\textwidth]{soc_heatmap.png}
% \caption{Annual heatmap of battery SoC for 2021 determined by health-aware RI. Upper panel shows hourly SoC with daily traded energy and mean SoC profile. Lower panel shows logarithmic distribution of SoC states. Clear daily periodicity is evident with higher charge levels in early morning and afternoon.}
% \label{fig:soc_heatmap}
% \end{figure}

\textbf{Observations:}
\begin{itemize}
\item \textbf{Daily periodicity}: Battery charges primarily during early morning (3-6 AM) and afternoon (2-4 PM) when prices are typically low
\item \textbf{Intermediate SoC states}: Counter to intuition, the battery spends significant time at intermediate SoC levels rather than purely full/empty states
\item \textbf{Increased activity in Q4}: Trading volume increases notably in October-December 2021, correlating with elevated gas prices and market volatility
\item \textbf{Complex trading patterns}: The LOB-based approach induces non-trivial SoC trajectories that differ from simplified bang-bang strategies
\end{itemize}

\subsection{Lifecycle Economics}
Table \ref{tab:lifecycle_economics} compares lifecycle metrics between degradation-aware and degradation-agnostic strategies.

% \begin{table}[h]
% \centering
% \caption{Lifecycle Economics: Health-Aware vs. Agnostic Strategies (2021)}
% \label{tab:lifecycle_economics}
% \resizebox{\textwidth}{!}{
% \begin{tabular}{lrrr}
% \toprule
% \textbf{Metric} & \textbf{No Degradation Cost} & \textbf{Linear ($\nu_{deg}=4$)} & \textbf{Health-Aware (16-seg)} \\
% \midrule
% Annual revenue (€) & 789,300 & 349,000 & 372,300 \\
% Annual life loss (\%) & 77.0 & 2.2 & 2.6 \\
% Prorated aging cost (€) & 2,887,500 & 81,300 & 96,300 \\
% Net annual profit (€) & -2,098,200 & 267,700 & 276,000 \\
% Battery life expectancy (yr) & 1.1 & 8.2 & 8.0 \\
% Cycles per day & 4.8 & 1.4 & 1.5 \\
% \bottomrule
% \end{tabular}
% }
% \end{table}

\textbf{Key Findings:}
\begin{enumerate}
\item The \textbf{degradation-agnostic} strategy maximizes short-term revenue but operates at a \textbf{massive loss} when accounting for battery replacement costs
\item Both degradation-aware strategies achieve \textbf{positive lifecycle profits} by intelligently managing cycle frequency and depth
\item The \textbf{health-aware piecewise model} achieves slightly higher profits (+3.1\%) than the simple linear model while providing more accurate degradation predictions
\item Battery life expectancy increases from 1.1 years (agnostic) to 8+ years (health-aware), approaching the 10-year shelf life limit
\item The health-aware strategy reduces cycling from 4.8 to 1.5 cycles/day while maintaining strong profitability
\end{enumerate}

\section{Discussion}
\subsection{Computational Efficiency}
Our DP formulation achieves 100-1000× speedup compared to exact MILP solutions through three key innovations:

\textbf{1. Decomposition by product:} Treating each delivery hour as a separate DP stage reduces the problem dimension dramatically. While the monolithic MILP has $O(T \cdot |\mathcal{O}|)$ decision variables, each DP stage involves at most $O(m \cdot |\mathcal{O}_t|)$ evaluations.

\textbf{2. Discretization:} Restricting actions to multiples of $v \cdot u$ limits the action space to $O((\bar{f} - \underline{f})/u)$ possibilities per stage, making enumeration feasible.

\textbf{3. Value function approximation:} Linear interpolation on grid $\mathcal{G}$ enables continuous state spaces while maintaining $O(m)$ storage and computation per stage.

The total computational complexity is $O(T \cdot m \cdot A)$ where $A$ is the number of feasible actions per stage, compared to the exponential complexity of branch-and-bound for MILP.

\subsection{Trade-offs and Limitations}
\textbf{Myopia:} The RI policy is fundamentally myopic, not anticipating future price changes. This leads to:
\begin{itemize}
\item Premature trading in illiquid markets far from delivery
\item Missed opportunities to wait for better prices
\item Inefficient use of storage flexibility
\end{itemize}
However, our parametric enhancement partially addresses this by penalizing illiquid trading.

\textbf{Perfect information assumption:} Our backtesting assumes perfect knowledge of how our orders would have been matched. In reality, algorithmic trading by multiple participants can affect market dynamics in ways we cannot capture from historical data alone.

\textbf{Linear degradation in segments:} While the piecewise approximation is accurate with sufficient segments, it remains an approximation. More sophisticated degradation models incorporating temperature, C-rate, and SoH could further improve accuracy.

\textbf{Single-market focus:} We consider only the continuous intraday market. Multi-market strategies coordinating day-ahead, intraday, and reserve markets could achieve higher profits but require more complex optimization.

\subsection{Future Research}
Promising extensions include: (1) stochastic DP incorporating price forecasts to reduce myopia, (2) multi-market optimization coordinating day-ahead, intraday, and reserve markets, (3) advanced degradation models incorporating temperature and C-rate effects, (4) deep reinforcement learning for non-myopic policies, and (5) adaptation to emerging markets such as the Indian Energy Exchange.

\section{Conclusion}
This paper presents a health-aware high-frequency trading framework for BESS in continuous intraday electricity markets. By integrating a piecewise linear battery degradation model \cite{xu2017factoring} with an efficient DP formulation inspired by Schaurecker et al. \cite{schaurecker2025maximizing}, we achieve computational speed sufficient for millisecond-pace trading while preserving long-term asset value.

Our DP formulation solves the intrinsic optimization 100-1000$\times$ faster than exact MILP methods, enabling reaction to every LOB update. The piecewise degradation model achieves $<2\%$ error with $J=16$ segments. Backtesting on 2021 German intraday data shows that high-frequency trading earns 58\% more than hourly optimization, health-aware strategies extend battery life from 1.1 to 8+ years, and parametric enhancements add a further 8.4\% improvement.

These results demonstrate that high-frequency trading is essential for maximizing BESS profitability, but must be balanced against battery health. Our framework achieves this balance, providing a practical approach to operating grid-scale storage in fast-moving electricity markets.

\section*{Acknowledgments}
The author thanks Dr. Debaprasad Sahu for supervision and guidance, and acknowledges the valuable contributions of the research papers by Xu et al. \cite{xu2017factoring} and Schaurecker et al. \cite{schaurecker2025maximizing} which form the theoretical foundation of this work.

\bibliographystyle{plainnat}
\begin{thebibliography}{99}
\bibitem{martinez2023dynamical}
Martínez-Barbeito, M., Gomila, D., and Colet, P. (2023).
Dynamical model for power grid frequency fluctuations: Application to islands with high penetration of wind generation.
\textit{IEEE Transactions on Sustainable Energy}, 14(3):1436--1445.

\bibitem{ullah2024comprehensive}
Ullah, F., Zhang, X., Khan, M., Mastoi, M.S., Munir, H.M., Flah, A., and Said, Y. (2024).
A comprehensive review of wind power integration and energy storage technologies for modern grid frequency regulation.
\textit{Heliyon}, 10(10):e30466.

\bibitem{braeuer2019battery}
Braeuer, F., Rominger, J., McKenna, R., and Fichtner, W. (2019).
Battery storage systems: An economic model-based analysis of parallel revenue streams and general implications for industry.
\textit{Applied Energy}, 239:1424--1440.

\bibitem{seifert2024coordinated}
Seifert, P.E., Kraft, E., Bakker, S., and Fleten, S.E. (2024).
Coordinated trading strategies for battery storage in reserve and spot markets.
\textit{arXiv preprint arXiv:2406.08390}.

\bibitem{xu2017factoring}
Xu, B., Zhao, J., Zheng, T., Litvinov, E., and Kirschen, D.S. (2017).
Factoring the cycle aging cost of batteries participating in electricity markets.
\textit{IEEE Transactions on Power Systems}, 33(2):2248--2259.

\bibitem{graf2022drives}
Gräf, D., Marschewski, J., Ibing, L., Huckebrink, D., Fiebrandt, M., Hanau, G., and Bertsch, V. (2022).
What drives capacity degradation in utility-scale battery energy storage systems?
\textit{Journal of Energy Storage}, 47:103533.

\bibitem{lohndorf2021gas}
Löhndorf, N. and Wozabal, D. (2021).
Gas storage valuation in incomplete markets.
\textit{European Journal of Operational Research}, 288(1):318--330.

\bibitem{semmelmann2025algorithm}
Semmelmann, L., Dresselhaus, J., Miskiw, K.K., Ludwig, J., and Weinhardt, C. (2025).
An algorithm for modelling rolling intrinsic battery trading on the continuous intraday market.
\textit{SIGENERGY Energy Inform. Rev.}, 4:163--174.

\bibitem{schaurecker2025maximizing}
Schaurecker, D., Wozabal, D., Löhndorf, N., and Staake, T. (2025).
Maximizing battery storage profits via high-frequency intraday trading.
\textit{European Journal of Operational Research} (forthcoming).
\end{thebibliography}
\end{document}